\chapter{Literature Review}\label{ch:literature}

\section{From lattice gases to hydrodynamics}
The FHP construction established that deterministic Boolean lattice gases on a triangular lattice can reproduce the Navier Stokes equations after coarse graining \cite{FHP1986}. Its importance for the present work is not merely historical. A lattice gas supplies an exactly specified microscopic map, local conservation laws, and finite state dynamics. Those features allow microscopic identities to be checked without floating point ambiguity while still connecting to continuum transport questions.

The separation between exact microdynamics and hydrodynamic closure is already present in the early literature. Poiseuille calculations were used to measure the effective shear viscosity of lattice gas automata and compare the macroscopic flow profile with kinetic predictions \cite{Kadanoff1987}. Green Kubo relations were subsequently derived for lattice gas cellular automata from microscopic conservation laws and correlation functions \cite{ErnstDufty1989}. These results make a useful methodological point: a transport coefficient can be represented through time correlations even when its approximate kinetic value is obtained by a separate molecular chaos closure.

For the present thesis, that distinction is decisive. The deterministic collision and streaming map can obey exact symmetry identities even if the correlated transport coefficient differs from its Boltzmann value. A proof at the level of a collision table therefore does not automatically extend to a proof about a zero frequency Green Kubo viscosity.

\section{Two dimensional correlations and limiting procedures}
Hydrodynamic modes generate algebraic correlation tails in low dimensional fluids. The general language of long wavelength fluctuations and correlation functions is developed in standard treatments such as Forster \cite{Forster1975} and in renormalization analyses of fluid fluctuations \cite{ForsterNelsonStephen1977}. For the two dimensional FHP cellular automaton specifically, Naitoh, Ernst, and Dufty found an intermediate $1/t$ regime in Green Kubo time correlation functions and a later self consistent asymptotic form proportional to $1/[t\sqrt{\ln t}]$ \cite{Naitoh1990}.

This literature rules out a casual instruction to remove long time tails as if they were numerical contamination. They are a physical part of the correlated dynamics. The correct response question must retain a regulator or specify the order of limits. One may use a finite lattice size $L$, a finite integration time $T$, a nonzero frequency $\omega$, or a running coarse graining scale $\ell$. A statement about circle survival can then be asked at fixed regulator before a limiting procedure is attempted.

The distinction also limits what is meant by viscosity in this thesis. The symbol $\eta$ may denote the published kinetic coefficient in \cref{ch:kinetic}, but a fully correlated object should carry its regulator or response prescription when ambiguity is possible. This convention prevents a finite Boltzmann coefficient from being silently identified with an infinite system Green Kubo limit.

\section{Odd viscosity}
Odd viscosity is the parity odd part of the viscosity response and is reactive rather than dissipative in the standard isotropic constitutive decomposition \cite{Avron1998}. In chiral active fluids, broken parity and time reversal can produce such a coefficient in a classical nonequilibrium setting \cite{Banerjee2017}. The broader literature emphasizes that odd transport is compatible with classical stochastic and kinetic models, not only quantum Hall systems \cite{FruchartReview2023}.

Several microscopic approaches are relevant. Fruchart \etal\ used non Hermitian kinetic theory to show how chiral collisions can generate Hall viscosity and odd thermal conductivity \cite{Fruchart2022}. Eren, Fruchart, and Vitelli developed a collisional odd fluid from Boltzmann kinetic theory and related hydrodynamic variables to spectral structure of a linear collision operator \cite{Eren2026}. These works motivate the use of complex eigenvalues and reduced response blocks, but they do not imply the particular anchored circle classified in this thesis.

Green Kubo methods have also been developed and tested for odd viscous fluids. Epstein and Mandadapu derived response relations for two dimensional nonequilibrium viscous fluids with broken time reversal symmetry \cite{Epstein2020}. Hargus \etal\ compared Green Kubo and nonequilibrium molecular dynamics routes for shear and odd viscosity in active fluids \cite{Hargus2020}. The existence of these methods supports the use of a complex correlated response $Z$, but a geometric relation between its real and imaginary parts still requires an additional statement.

\section{Kubo formulas and Ward identities}
Kubo's linear response framework expresses transport through equilibrium or stationary correlation functions \cite{Kubo1957}. For viscosity, care is needed in defining stress response and contact terms. Bradlyn, Goldstein, and Read derived viscosity Kubo formulas from strain generators and analyzed Ward identities that follow from symmetries and, under additional assumptions, relate stress response to conductivity \cite{Bradlyn2012}.

This provides the standard against which any claimed Ward identity should be judged. A Ward identity is not merely a compact algebraic relation. It is tied to a symmetry or conservation law. The response circle identity in this thesis arises from complex inversion and an anchored geometric equation. It becomes Ward like only if a microscopic principle can be found that forces the dissipative inverse response to remain chirality independent. No such principle is established in the present work.

The same caution applies to Kramers Kronig language. Causal analyticity links the real and imaginary parts of a frequency dependent response. The CRTC problem instead follows a control parameter $p$ at fixed response prescription. A geometric relation in the chirality parameter plane is therefore not a nonlinear Kramers Kronig relation. The mathematical resemblance between dissipative and reactive coordinates is not a substitute for the analyticity variable required by the Kramers Kronig construction.

\section{The 2026 chiral FHP model}
Allocca \etal\ introduced parity breaking two body collision rules and a systematic velocity rotation into the FHP framework, derived hydrodynamic transport coefficients, and compared them with Poiseuille flow simulations \cite{Allocca2026}. Their zero field shear sector has a complex relaxation eigenvalue whose real part is independent of the chirality probability $p$ within the molecular chaos calculation. The viscosity pair is an inverse image of that vertical eigenvalue locus.

This observation contains the seed of the response geometry studied here. At fixed density, a vertical line in inverse response space is mapped by complex inversion to a circle through the origin. The circle is therefore a geometric restatement of a specific spectral flatness property in the kinetic model. Once correlations are admitted, the useful question is whether the corresponding real part remains flat after the effective response is computed.

The present thesis relies on the published kinetic formulas as an external premise in \cref{ch:kinetic}. It then reconstructs the algebra independently and checks it exactly over finite rational domains. This is intentionally narrower than claiming the kinetic derivation as original work.

\section{Classical response circles and novelty boundary}
Circular arcs are not new objects in complex response theory. Complex inversion maps generalized circles and lines to generalized circles and lines, and single relaxation response models have long produced semicircular plots. Therefore the novelty claim in this thesis is not that a complex transport coefficient can trace a circle.

The contribution is instead the following combination. First, the chiral FHP kinetic coefficients are placed into a complex inverse response representation with an exact executable reconstruction. Second, the anchored circle is characterized by a necessary and sufficient condition on the measured inverse response that does not require a fixed kinetic prefactor. Third, symmetry is separated from rigidity by an exact commutant calculation. Fourth, a positive dissipation correlated Green Kubo family is constructed that satisfies the symmetry form but violates the circle. Fifth, the role of spectral information is qualified by distinguishing an identified reduced shear block from a bare unordered full spectrum.

This boundary matters because each element answers a different question: prior work supplies the hydrodynamic and transport setting; complex geometry supplies a mathematical mechanism; the present classification identifies the exact invariant that must survive renormalization and provides an explicit witness when it does not.

\section{Chapter conclusion}
The literature supports three premises for the analysis that follows. Deterministic lattice gases can possess exact microscopic structure while their transport is obtained by coarse graining. Two dimensional correlation tails make the fully correlated transport problem nontrivial. Odd viscosity and complex shear response are established concepts, but neither sixfold symmetry nor the existence of a Hall viscosity is known to impose the CRTC circle. These premises motivate an evidence hierarchy that keeps microscopic theorems, kinetic corollaries, and correlated transport conjectures separate.
